\section{Related Work}
\label{sec:related_work}

\textbf{Static Weight-Space Merging:} Model merging began with simple weight averaging schemes, such as Model Soups \cite{Wortsman2022}, which averages fine-tuned weights from the same pre-trained initialization to improve generalization. Task Arithmetic \cite{Ilharco2022} extended this by adding and subtracting task vectors (fine-tuned weights minus pre-trained base weights). To address parameter interference and sign conflicts, TIES-Merging \cite{Yadav2023} and DARE \cite{Yu2023} introduced pruning, sign consensus, and rescaling. Other static methods like RegMean \cite{Jin2022} align activations in closed-form, and ZipIt! \cite{Stoica2023} merges models trained from different initializations by aligning features. However, static fusions are fundamentally limited to finding a single consensus model, which often degrades performance when merging highly distinct or orthogonal tasks.

\textbf{Dynamic Routing in Model Merging:} To bypass the constraints of static merging, recent research has explored dynamic merging, adapting Mixture-of-Experts (MoE) routing concepts \cite{Shazeer2017, Fedus2021, Riquelme2021} to weight-space blending. Instead of running multiple separate feed-forward expert paths (which increases inference compute), dynamic merging routes input representations at test-time to compute sample-specific blending coefficients, merging the weights on the fly so that inference is performed on a single fused model \cite{PredecessorT4S6}. This concept of projecting penultimate feature representations onto pre-trained expert classification weights to identify task alignment is closely connected to \emph{Prototypical Networks} \cite{Snell2017} and zero-shot classification in metric learning, where sample similarities to semantic class prototypes are leveraged for classification without explicit optimization. 

\textbf{Over-Parameterization and Deconstruction:} Recent conference contributions have introduced highly elaborate routing mechanisms. Most notably, Quantum Wavefunction Superposition Merging (QWS-Merge) \cite{PredecessorT4S10} uses quantum wave-inspired metaphors to generate merging coefficients, claiming unique robustness. However, subsequent rigorous investigations \cite{PredecessorT5S4, PredecessorT5S5} deconstructed these wave-inspired formulations, proving they suffer from severe transductive overfitting \cite{PredecessorT2S1} on OOD data (e.g., SVHN). They demonstrated that the wave-like phase activation functions behave in an unstable manner, whereas a standard classical linear router with simple $L_2$ weight decay regularization achieves superior generalizability and robustness. 

Our work builds directly on these deconstructions. We take a radical, minimalist step forward: instead of optimizing a regularized classical router, we eliminate trainable parameters and calibration-tuning completely.

\textbf{The Robustness-Accuracy Illusion:} Layer-wise dynamic routers, such as the L3-Softmax Router \cite{PredecessorT5S5}, use Softmax normalization at every layer to enforce simplex constraints. While they exhibit relative stability under varying batch configurations, previous audits \cite{PredecessorT5S5} unmasked this as a \emph{Robustness-Accuracy Illusion}. Simplex normalization restricts the router's expressivity, forcing the dynamic coefficients to hover near a uniform average ($\approx 0.25$ for four tasks) regardless of the input. This flat routing prevents the model from specializing, leading to mediocre, uniform-like performance. 

In contrast, our proposed PFSR + MBH framework achieves high task specialization (achieving a Joint Mean of $75.00\%$, close to the expert ceiling of $79.80\%$) while remaining completely immune to heterogeneity collapse, proving that solving stream issues at the data level is strictly superior to over-constraining the model architecture.

\textbf{High-Performance Serving and Request Batching:} Our proposed stream partitioning and micro-batch assembly in MBH is conceptually connected to high-performance serving and request batching frameworks in systems-ML, such as Orca \cite{Yu2022} and vLLM \cite{Kwon2023}. These systems leverage dynamic request partitioning and continuous batching to optimize hardware utilization and throughput under heterogeneous generation workloads. However, while Orca and vLLM focus on scheduling and memory management (such as PagedAttention) for single or multiple standalone models, MBH applies dynamic partitioning directly to weight-space model merging, grouping heterogeneous input streams into task-homogeneous micro-batches to eliminate representation interference and weight-averaging collapse. This establishes a unique bridge between systems-level request scheduling and parameter-level dynamic model merging.
